\section{问题重述}
\label{sec:1}

\subsection{问题背景}
\label{sec:1.1}

大语言模型的预训练能力随模型参数量 $N$（个）与训练 token 数 $D$（token）增长，经典标度律以验证集交叉熵损失 $L_0(N,D)=E+AN^{-\alpha}+BD^{-\beta}$ 刻画这一过程\cite{hoffmann,bahri}。按本题采用的近似，基础训练开销约为 $C\approx 6ND$（FLOPs），有限预算下必须权衡模型规模与数据量。仅用 $N,D,C$ 不足以刻画真实训练：语料质量与领域组成会显著改变同等投入的产出。为此，需要先建立数据质量 $Q$ 的可比口径与领域配比 $\bp$ 的损失关系，再将二者接入经典标度律，形成统一的 $L(N,D,Q,\bp)$，最终落到算力约束下的资源配置与能力演进判断。附件 A 提供质量信号与配比实验，附件 B 提供标度律数据，附件 C 提供模型、算力与能力评测信息。

\subsection{问题描述}
\label{sec:1.2}

本题依次要求完成四项数学任务，前一问的输出作为后一问的输入。

\textbf{问题一：数据质量与领域配比。} 建模对象为文档、领域与混合语料的质量，以及 17 个训练领域配比对 13 个验证域损失的影响；主要约束是质量信号方向与量纲不一、11 个配比域缺少直接质量映射，且配比份额非负、总和为 1。本文用 A1（51,230 篇）拟合文档质量口径与全部预处理参数，冻结后在扩展集 A2（arxiv 17,523 篇）、A3（github 203,752 篇）上检验评分与诊断的稳定性；用 A16 映射与 A18 文本特征把质量扩展到 17 域；A4/A5 按 \texttt{index} 连接作 1M 配比训练，A6/A7 作 1M 留出，A8/A9 作 60M 跨尺度检验与尺度修正标定，A10/A11（各 64 行）作 1B 独立检验，A12–A15 用于估算/外推情景讨论。

\textbf{问题二：广义标度律。} 建模对象为损失对 $N,D,Q,\bp$ 的依赖；若采用理想质量与参照配比下退回经典式的约束，须说明依据；且 A、B 两组损失的评测管线与质量尺度不可直接等同。本文以 B1 估计经典基线，用 B6（360 点，半合成）估计质量项并选择模型，用 B7 相对 B6 新增的 90 点（$Q=0.5,0.7$，覆盖全部 45 个 $(N,D)$ 格）作留出；B2、B3、B4、B5、B8、B9、B10 分别用于半合成、插值、跨族、文献、方向冲突与估算对照，B11/B12 用于来源审计。问题一的域质量与配比关系经固定接口引入。

\textbf{问题三：算力约束下的资源配置。} 建模对象为给定预算 $C$（FLOPs）与上下文长度 $L_{ctx}$（token）下的 $N,D,Q$；约束包含训练、提质与长上下文三项成本，质量下界由问题一推荐配比确定，主方案固定该配比。本问读取问题一的质量与配比接口、问题二冻结的标度律参数，不在本问重新拟合；C7 提供上下文长度的可行档位。

\textbf{问题四：Benchmark 检验与能力演进。} 建模对象为开放权重模型的能力前沿及其规模与非规模来源；要求能力、时间、模型类型与开放权重口径保持一致，且损失与 Benchmark 得分不在同一尺度。C1/C2 提供榜单能力与发布日期，C4 补充规模与算力，C8 用于逐任务聚合与饱和度检查，C6 分层拟合损失到得分的桥接，C3 仅作混合口径的敏感性对照。
